% !TeX root = ../../Book3.tex
\OptionalSection{曲率方法}{b3:sec:1406}

本选读节说明负曲率与重标度如何组织更广的全纯刚性。首先保留统一曲率常数的准确范围，再以紧目标讨论整个曲线与正规性的联系。

\subsection{\texorpdfstring{\(\mathbb C\)}{C} 上不存在负曲率度量}
\label{b3:subsec:140601}

\begin{corollary}\label{b3:cor:14-no-uniform-negative-plane}
全平面不存在正 \(C^2\) 共形密度满足 \(K\le-\kappa<0\)。
\end{corollary}
\begin{proof}
对任意中心 \(a\) 的半径 \(R\) 圆盘使用 Ahlfors--Schwarz，在中心得到 \(\rho(a)\le2/(\sqrt\kappa R)\)。令 \(R\rightarrow\infty\)，得到 \(\rho(a)=0\)，违反正性。
\end{proof}
这里只排除统一负上界，不排除处处负曲率。\(\rho=e^{|z|^2}\) 给 \(K=-4e^{-2|z|^2}<0\)，且度量完备，提供准确反例。标题中的“负曲率”应按前述量化陈述理解，不能删除 \(\kappa\) 而使定理变成错误命题。

\subsection{Liouville 型结论}
\label{b3:subsec:140602}

若目标有统一负曲率密度，任意整平面全纯映射到它都常值。把映射限制在以任意源点为中心的大圆盘，导数比较给 \(\rho(f(a))|f'(a)|\le2/(\sqrt\kappa R)\)；令 \(R\rightarrow\infty\)，目标密度正使 \(f'(a)=0\)。这是一种几何 Liouville 结论，目标未必在某个坐标中有界。

\ProseParagraphBreak
对于双曲曲面，Poincaré 密度曲率 \(-1\)，所以整曲线均常值；也可直接提升到圆盘使用 Liouville。两个证明分别突出曲率与覆盖，所需前提等价于已建立的双曲模型。

\subsection{Brody 型思想}
\label{b3:subsec:140603}

重标度还适用于紧 Riemann 曲面目标的任一光滑 Hermitian 度量。若圆盘映射列导数在内点附近无共同界，按本章最大化“到边界距离乘导数”的方法选中心与尺度，得到越来越大圆盘上导数有共同界、中心导数为1的映射列。目标紧给取值紧性，等度连续性给全平面极限；在局部复坐标中用 Cauchy 公式使极限全纯且中心导数非零。因此出现非恒常整曲线。

\ProseParagraphBreak
反向，若有非恒常整曲线 \(g:\mathbb C\rightarrow X\)，在某个非零导数点平移后考察 \(g(nz)\)，其中心导数无界，故圆盘映射族不正规。这是曲面情形的 Brody 重标度思想。紧性不可省略，否则像可能逃向目标外部而没有全平面极限；球面 Zalcman 证明可工作，正因为球面本身紧。

\subsection{双曲性理论的进一步讨论}
\label{b3:subsec:140604}

紧球面有非恒常整曲线，例如 \(z\mapsto z\) 视为球面值映射；复环面也有投影 \(\mathbb C\rightarrow\mathbb C/\Lambda\)。紧的圆盘型曲面则没有这样的曲线。因此全纯曲线刚性识别了紧曲面三种覆盖模型中的双曲情形。

\ProseParagraphBreak
更高维中，负曲率、Kobayashi 距离和整个曲线之间的关系需另立定义与定理，不能只把一维 Gaussian 曲率公式改几个指标后宣称全部成立。本卷以后讨论双曲性时保留一维范围；本节只提供从导数比较、紧性到全局刚性的证明思路。

\ProseParagraphBreak
本章有两种从局部估计取得全局结论的方法。曲率比较在越来越大的圆盘上控制拉回密度，极限使整映射的导数消失；重缩放则从坏的局部导数中提取一个非平凡整映射，再以省略值或目标双曲性排除它。两者所用的几何尺度不同，但都必须核对新定义域是否真的趋于整个平面，且归一化是否保留非平凡性。

\ProseParagraphBreak
在球面和环面中，非恒定整曲线确实存在：前者可取 \(z\mapsto z\) 作为球面值映射，后者取商投影 \(\mathbb C\to\mathbb C/\Lambda\)。圆盘覆盖的曲面则没有，源平面单连通使它提升到圆盘，再由 Liouville 为常数。这个区别是在本书已经标明的统一化输入上得到的完整推论；更一般复流形的 Brody 定理需要额外紧致性和几何论证，不在这里扩大声称。
